\subsection{Direct local enclosure calculus}
\label{subsec:direct-local-calculus}

Let $\mathsf R$ denote counterclockwise rotation through $2\pi/3$, and put
$h=\sqrt3/2$.  For a nonempty compact set $K$, let
\[
 h_K(n)=\max_{x\in K}\langle x,n\rangle .
\]

\begin{proposition}[Equilateral enclosure gauge]
\label{prop:new-enclosure-gauge}
\label{prop:universal-enclosure-gauge}
The least side length of a closed equilateral triangle containing $K$ is
\[
 \Lambda(K)=\frac1h\min_{\lVert n\rVert=1}
 \sum_{j=0}^2h_K(\mathsf R^jn).
 \tag{6.1}
\]
If $K$ is a polygon, a minimizing normal can be chosen perpendicular to a
hull edge.
\end{proposition}

\begin{proof}
For fixed outward normals $n,\mathsf Rn,\mathsf R^2n$, the least containing
triangle has those three support numbers, and its side is their sum divided
by $h$.  On an open angular cell with fixed support points, the support sum is
$A\cos\theta+B\sin\theta$ and has second derivative equal to its negative.
Thus an interior critical point is a maximum, and a minimum occurs at a
support tie, equivalently at a hull-edge normal.
\end{proof}

Take unit vectors $u,v$ meeting at $120^\circ$, so that $u+v$ is the inward
radial direction, and put
\[
 K(a,b,c)=\conv\{0,au,bv,c(u+v)\}.
\]
The demand triple $(a,b,c)$ is \emph{admissible} when $K(a,b,c)$ lies in a
closed unit equilateral triangle.

\begin{proposition}[Exact local admissible set]
\label{prop:new-exact-local-set}
\label{prop:exact-admissible-set}
Put
\[
 s=a+b,\quad d=a^2+ab+b^2,\quad
 m=\min\{a,b\},\quad M=\max\{a,b\},
\]
\[
 q=s^4-s^2+ab,
\]
and
\begin{align*}
 F_L&=c^4-c^2+mc-m^2,\\
 F_T&=(s^2-1)c^2+Mc-M^2,\\
 F_S&=(m^2-1)c^2+(2mM^2+M)c+M^4-M^2.
\end{align*}
Then the admissible set is the union of
\begin{align*}
 \mathcal A_L&=\{d\le1,\,s\le1,\,q\le0,\,F_L\le0\},\\
 \mathcal A_T&=\{d\le1,\,s\le1,\,q\ge0,\,F_T\le0,\,c\le2M\},\\
 \mathcal A_S&=\{d\le1,\,s>1,\,F_S\le0,\,c\le1/2\}.
\end{align*}
Consequently the maximal admissible radial demand is
\[
 c_{\max}(a,b)=
 \begin{cases}
  1,&a=b=0,\\
  C_L(m),&s\le1,\,q\le0,
  \quad (a,b)\ne(0,0),\\
  \dfrac{2M}{1+\sqrt{4s^2-3}},&s\le1,\,q>0,\\[3mm]
  \dfrac{M(2mM+1)-M\sqrt{4d-3}}{2(1-m^2)},&s>1,
 \end{cases}
 \label{eq:radial-envelope}
 \tag{6.2}
\]
where $C_L(m)$ is the selected root in $[\sqrt3/2,1]$ of
\[
 z^4-z^2+mz-m^2=0.
\]
The admissible set is coordinatewise down-closed.
\end{proposition}

\begin{proof}
If $cs\le ab$, the radial point lies in $\conv\{0,au,bv\}$ and the least
side is the diameter $\sqrt d$.  Otherwise the convex hull has four edges.
Proposition~\ref{prop:new-enclosure-gauge} reduces the minimum to the four
edge normals.  The corresponding side lengths are
\begin{align*}
 L_{0A}&=b+\max\{a,c\},&
 L_{B0}&=a+\max\{b,c\},\\
 L_{AC}&=
 \begin{cases}
 (a^2+bc)/\sqrt{a^2-ac+c^2},&a\ge c,\\
 c(a+b)/\sqrt{a^2-ac+c^2},&a<c\le a+b,\\
 c^2/\sqrt{a^2-ac+c^2},&c>a+b,
 \end{cases}
\end{align*}
and $L_{CB}$ is obtained by interchanging $a,b$.  Squaring their positive
unit-frontier equations gives $F_L,F_T,F_S$.  The inactive support conditions
give respectively $q\le0,q\ge0,s>1$.  In the $T$ cell the selector
$c\le2M$ removes the spurious large root; in the $S$ cell $c\le1/2$ selects
the geometric component.  Solving the selected equations gives (6.2).
Convexity of a containing triangle proves down-closedness.
\end{proof}

The zero-radial diameter envelope is
\[
 M_0(x)=\frac{-x+\sqrt{4-3x^2}}2.
 \tag{6.3}
\]
For
\[
 0<e<1-\frac{\sqrt3}{2},
\]
define the selected low root
\[
 \epsilon(e)=\frac{1-e}{2}
 \left(1-\sqrt{4(1-e)^2-3}\right).
 \tag{6.4}
\]
It satisfies
\[
 e<\epsilon(e),\qquad
 \epsilon(e)<\frac{2e}{1-2e},
 \tag{6.5}
\]
and, for $e\le1/8$,
\[
 \epsilon(e)\le2e+5e^2.
 \tag{6.6}
\]
These follow by substituting the proposed upper bounds in the selected
quadratic and checking its sign.

\begin{lemma}[Direct high-radial threshold]
\label{lem:new-direct-threshold}
Let a nonsupercritical V triangle have incoming boundary reach $A$ and
outgoing boundary reach $B$.  If
\[
 C\ge1-e,
 \qquad
 A>\epsilon(e),
\]
then
\[
 \boxed{B\le\epsilon(e).}
 \tag{6.7}
\]
\end{lemma}

\begin{proof}
The condition $C>1/2$ forces $A+B\le1$.  In the exact cells of
Proposition~\ref{prop:new-exact-local-set}, the only selected component with
$A>\epsilon(e)$ has $B\le\epsilon(e)$; equivalently the two roots of the
selected $T$-cell quadratic separate the feasible fiber.  The selector
$c\le2\max\{A,B\}$ removes the formal high root.
\end{proof}

For $0\le c\le1/2$, the strict-supercritical outgoing supremum is
\[
 M_c^{\rm sup}
 =\frac{c+\sqrt{c^2-8c+4}}2.
 \tag{6.8}
\]
It is not attained.  Thus an actual strict-supercritical V triangle with
own-radial reach at least $c$ has outgoing boundary reach strictly below
$M_c^{\rm sup}$.

\begin{lemma}[Selected-arc chord bounds]
\label{lem:new-selected-chords}
Let $0<e<1/2$.  On the selected $Q_+$ component for radial demand $1-e$,
let $p$ be the incoming boundary reach and let $q$ be the following incoming
reach forced across a center-free edge.  Then
\[
 q-p=1-\sqrt{1-x+x^2},
\]
where
\[
 q=e+(1-e)x,
 \qquad
 p=e+(1-e)x-1+\sqrt{1-x+x^2}.
 \tag{6.9}
\]
The selected arc is increasing and concave.  In particular, on every interval
used below,
\[
 q\ge p+\lambda(p-e)
 \tag{6.10}
\]
for the chord slope $\lambda$ joining the appropriate selected endpoints.
The CE1 applications may take
\[
 \lambda>1-4e
 \quad(e=\alpha),
 \qquad
 \lambda>1-5e
 \quad(e=\alpha/R).
 \tag{6.11}
\]
\end{lemma}

\begin{proof}
The selected $T$-cell equation reduces, after putting
$x=(q-e)/(1-e)$, to
\[
 x(1-x)=(q-p)(2-q+p).
\]
The selected component takes the smaller root, giving (6.9).  Strict
concavity follows by differentiation.  A concave graph lies above its chords.
For the first slope, (6.5) gives
\[
 \frac{e}{\epsilon(e)-e}>1-4e.
\]
For the second, use (6.6) when $e\le1/8$ and
$\epsilon(e)<(1-e)/2$ otherwise; the resulting slope is greater than
$1-5e$ throughout $0<e<1-\sqrt3/2$.
\end{proof}
